% Options for packages loaded elsewhere
\PassOptionsToPackage{unicode}{hyperref}
\PassOptionsToPackage{hyphens}{url}
\documentclass[
  a4paper,
]{article}
\usepackage{xcolor}
\usepackage[margin=25mm]{geometry}
\usepackage{amsmath,amssymb}
\setcounter{secnumdepth}{-\maxdimen} % remove section numbering
\usepackage{iftex}
\ifPDFTeX
  \usepackage[T1]{fontenc}
  \usepackage[utf8]{inputenc}
  \usepackage{textcomp} % provide euro and other symbols
\else % if luatex or xetex
  \usepackage{unicode-math} % this also loads fontspec
  \defaultfontfeatures{Scale=MatchLowercase}
  \defaultfontfeatures[\rmfamily]{Ligatures=TeX,Scale=1}
\fi
\usepackage{lmodern}
\ifPDFTeX\else
  % xetex/luatex font selection
\fi
% Use upquote if available, for straight quotes in verbatim environments
\IfFileExists{upquote.sty}{\usepackage{upquote}}{}
\IfFileExists{microtype.sty}{% use microtype if available
  \usepackage[]{microtype}
  \UseMicrotypeSet[protrusion]{basicmath} % disable protrusion for tt fonts
}{}
\makeatletter
\@ifundefined{KOMAClassName}{% if non-KOMA class
  \IfFileExists{parskip.sty}{%
    \usepackage{parskip}
  }{% else
    \setlength{\parindent}{0pt}
    \setlength{\parskip}{6pt plus 2pt minus 1pt}}
}{% if KOMA class
  \KOMAoptions{parskip=half}}
\makeatother
\usepackage{longtable,booktabs,array}
\newcounter{none} % for unnumbered tables
\usepackage{calc} % for calculating minipage widths
% Correct order of tables after \paragraph or \subparagraph
\usepackage{etoolbox}
\makeatletter
\patchcmd\longtable{\par}{\if@noskipsec\mbox{}\fi\par}{}{}
\makeatother
% Allow footnotes in longtable head/foot
\IfFileExists{footnotehyper.sty}{\usepackage{footnotehyper}}{\usepackage{footnote}}
\makesavenoteenv{longtable}
\usepackage{graphicx}
\makeatletter
\newsavebox\pandoc@box
\newcommand*\pandocbounded[1]{% scales image to fit in text height/width
  \sbox\pandoc@box{#1}%
  \Gscale@div\@tempa{\textheight}{\dimexpr\ht\pandoc@box+\dp\pandoc@box\relax}%
  \Gscale@div\@tempb{\linewidth}{\wd\pandoc@box}%
  \ifdim\@tempb\p@<\@tempa\p@\let\@tempa\@tempb\fi% select the smaller of both
  \ifdim\@tempa\p@<\p@\scalebox{\@tempa}{\usebox\pandoc@box}%
  \else\usebox{\pandoc@box}%
  \fi%
}
% Set default figure placement to htbp
\def\fps@figure{htbp}
\makeatother
\setlength{\emergencystretch}{3em} % prevent overfull lines
\providecommand{\tightlist}{%
  \setlength{\itemsep}{0pt}\setlength{\parskip}{0pt}}
\usepackage{bookmark}
\IfFileExists{xurl.sty}{\usepackage{xurl}}{} % add URL line breaks if available
\urlstyle{same}
\hypersetup{
  hidelinks,
  pdfcreator={LaTeX via pandoc}}

\author{}
\date{}

\begin{document}

\section{Causal Separation of the Time Structure of Three-Direction
Formation in Closed N-Body Relational Wave Systems via Two-Stage Seed
Removal}\label{causal-separation-of-the-time-structure-of-three-direction-formation-in-closed-n-body-relational-wave-systems-via-two-stage-seed-removal}

\textbf{Author:} Noriaki Kihara (WF System Co., Ltd.) \textbf{Date:}
2026-07-27 \textbf{Version DOI:}
\href{https://doi.org/10.5281/zenodo.21614403}{10.5281/zenodo.21614403}
\textbf{Concept DOI:}
\href{https://doi.org/10.5281/zenodo.21614402}{10.5281/zenodo.21614402}

\begin{center}\rule{0.5\linewidth}{0.5pt}\end{center}

\subsection{Claims}\label{claims}

Paper 7 discovered that a zero-square-closed relational system undergoes
geometric rapid expansion after a long low-change region and, through
that process, spontaneously forms a metastable three-directional
structure that did not exist initially. This is the central experimental
fact of this research. The present paper does not question that
three-direction formation. The purpose of Paper 8 is to examine whether
another bottom, repeating the same formation process one level below,
exists inside the seemingly flat region before the rapid expansion; to
verify, at the same time, whether the metastable state after the rapid
expansion can produce a further rapid-expansion cycle; and to separate
and clarify the time structure leading to three-direction formation and
the subsequent long-time evolution.

The verification found that, at least for N=5 within double-precision
computation, no separate bottom equipped with independent rapid
expansion, arrest, and metastabilization exists inside the initial flat
region. Even with the explicit seeds removed, the system starts the same
geometric growth from an extremely small state and passes the major
amplitude levels at almost the same times. The seemingly flat interval
is therefore not a stationary phase but a long latent process continuing
into the later rapid expansion.

From this latent stage onward, rank\_q=4 in the matrix diagnostics
already appears intermittently, persists during the geometric rapid
expansion, and becomes a clear and sustained structure in the metastable
region after the expansion. After the rapid expansion, new directional
occupations grow to finite magnitudes and are retained; that a
three-directional structure was actually formed is an undeniable
observational fact.

At the same time, detailed direction-lineage tracking revealed that the
small direction subspace before the rapid expansion was not simply
amplified while fixed to the completed three directions. The direction
subspace actually rotates, mixes, and reorganizes during the rapid
expansion. This does not weaken three-direction formation; rather, it is
a new finding showing that the three directions did not exist as
pre-fixed axes but were dynamically selected and formed inside the
rapid-expansion process.

This paper further verified the possibility that the metastable state
formed after the rapid expansion is an intermediate phase preparing the
next rapid expansion. In Paper 7, a second seed was added to this
metastable state and the subsequent time evolution was observed; within
the observation range in the time direction used since Paper 6, no
second rapid expansion was found. In this paper, that second seed itself
was removed and the natural long-time evolution of the metastable state
was observed. As a result, the metastable state settles into a steady
oscillatory state within roughly one period, and even when the
observation range along the time axis is extended, no sign connecting to
a second geometric rapid expansion was observed.

Therefore, at least within the range observed here, the metastable state
after the rapid expansion was not confirmed to be a relay point that
automatically produces the next rapid expansion. The metastable state
persists for a long time, but no clear growth series heading toward a
second rapid expansion has been found inside it. This result indicates
that the first rapid expansion may not be a universally repeating cycle,
and at the same time suggests that the state after three-direction
formation possesses an independent stabilization mechanism.

The phenomenon clarified by this paper is a single sequence of time
evolution: a long latent process, geometric rapid expansion, direction
reorganization during the expansion, spontaneous formation of a
metastable three-directional structure, and long-time metastabilization
without a second rapid expansion. What remains unexplained is not the
existence of the phenomenon. It is the formation and arrest mechanisms:
why the latency time is needed, why the geometric rapid expansion
begins, why the reorganization during the expansion selects three
directions, and why the formed metastable state settles down instead of
proceeding to the next rapid expansion.

An important point about the initial state of this experimental system
must be stated explicitly. The experiments initialize the N(N-1)/2
relational waves obtained from N entities by make\_parent, and this
initialization presupposes the zero-square closure condition, the first
axiom of this paper series. The initial state is therefore not an
unconditional placement of arbitrary continuous values; it is
constructed from the outset as a closed relational system.

As examined in a separate note, this closure condition inherently
contains quantization. Closable states do not exist continuously without
restriction; the closure condition restricts states to specific phase
relations. From this viewpoint, quantization is also demanded of the
complex phase, and the phase may be discretized in units of the \(\pi\)
period. The implementation in this paper does not explicitly impose such
a phase-quantization rule from outside.

Nevertheless, the many relational waves generated by make\_parent
consequently collapsed, in the initial state, onto two orthogonal
directions and formed a stationary two-directional state. The
experimenter did not specify two axes as an initial condition. What was
given was the relational system and the closure condition; the
two-directional structure appeared spontaneously as their consequence.

Therefore, the three-direction formation observed in this paper is not
an experiment that added a third axis to two artificially placed axes.
It is a phenomenon in which a two-directional state, spontaneously
established from zero-square closure, transitions to a metastable
three-directional structure through the long latent process, the
geometric rapid expansion, and the direction reorganization during the
expansion. Whether this initial collapse onto two directions is the
direct cause of three-direction formation has not yet been identified in
this paper. However, that the initial two directions were not introduced
from outside but arose from inside a relational system satisfying the
first axiom is an indispensable experimental fact for interpreting the
phenomenon.

This paper series has also confirmed, in separate experimental systems,
that wave interactions have at least two distinct types. One is a
bosonic linear interaction in which waves exchange and mix while
remaining superposable; the other is a fermionic interaction that
produces nonlinear scattering with reflection and transmission. In the
latter interaction, localization is exchanged between an extended wave
and a localized wave, and a metastable intermediate state is formed.

However, in the experiments so far, this fermionic wave or nonlinear
scattering rule did not arise spontaneously from inside the closed
relational system. It is explicitly given from outside by the
experimenter as the interaction type. Therefore, at the present stage it
cannot be said that the geometric rapid expansion and three-direction
formation observed in this paper can be explained by spontaneous
emergence of the fermionic interaction.

This paper also searched for precursors of the possibility that, in the
process of transition from the initial two-directional state to the
metastable three-directional structure, a distinction between linear
waves and nonlinearly reflecting-scattering waves emerges from inside
--- that is, the possibility that the separation into bosonic and
fermionic states is related to the onset of the rapid expansion or of
direction formation. However, within the time series, direction
occupations, rank changes, and metastabilization process observed here,
no clear sign was found by which such a distinction could be judged to
have formed spontaneously.

Therefore, the possibility that the separation of bosonic and fermionic
interactions is the cause of the geometric rapid expansion or of
three-direction formation is not excluded, but neither is it supported
by this experiment. What was confirmed is that, without installing a
fermionic scattering rule, the two-directional state established from
zero-square closure transitions to a metastable three-directional
structure through spontaneous rapid expansion and direction
reorganization. Whether the statistical distinction of waves appears
afterwards as a result of this formation process, or requires another
degree of freedom not yet observable in this experiment, remains
unresolved.

\begin{center}\rule{0.5\linewidth}{0.5pt}\end{center}

\subsection{1. System (definitions and
axioms)}\label{system-definitions-and-axioms}

\textbf{Axiom 1 (nontrivial zero-square-sum closure)}
\(\sum_e Z_e^2=0,\ Z\neq0\). \textbf{Axiom 2 (finite recurrence)}
\(U^n=I\). Both axioms follow the basic axiom system {[}1{]}. The system
definition is identical to Paper 7 {[}2{]}; the dynamical code reuses
the Paper 7 originals fixed by SHA-256 (§10).

A complex amplitude \(Z_e\) is placed on each edge \(e\) of the complete
graph \(K_N\) (\(M=N(N-1)/2\)). The state is closed on

\[Z\in\mathbb C^M,\qquad \lVert Z\rVert^2=1,\qquad Z^{\mathsf T}Z=0.\]

The generator \(K\) is a real antisymmetric matrix depending only on the
edge phases \(\theta_e=\arg Z_e\), coupling vertex-sharing edges through
the sine of phase differences:

\[(KZ)_e=\cos\theta_e\!\sum_{e'\sim e}\!\sin\theta_{e'}Z_{e'}-\sin\theta_e\!\sum_{e'\sim e}\!\cos\theta_{e'}Z_{e'}.\]

The update is the Cayley transform after spectral-norm normalization,
\(Z\leftarrow(I-\gamma\tilde K)^{-1}(I+\gamma\tilde K)Z\)
(\(\tilde K=K/\sigma_{\max}\), \(\gamma=\tan(\pi/144)\)). Since \(K\) is
antisymmetric this is unitary and conserves \(\lVert Z\rVert^2\)
exactly.

The parent state \(v\) is the self-consistent circularly polarized
eigenmode (all edges rotating rigidly at one frequency,
\(v^{\mathsf T}v=0\)). The parent plane is
\(P_1=\mathrm{span}\{\mathrm{Re}\,v,\mathrm{Im}\,v\}\). The eigenmode
residual of the double-precision parent is \(2.140\times10^{-13}\) for
N=5.

\subsubsection{1.1 Observables}\label{observables}

\begin{itemize}
\tightlist
\item
  \textbf{Splitting fraction}
  \(f(t)=\lVert Z-\Pi_{P_1}Z\rVert^2/\lVert Z\rVert^2\): total
  occupation outside the parent plane.
\item
  \textbf{crossing}: first passage of the existing Paper 7 threshold
  \(f>0.05\). A reference coordinate, not an event definition.
\item
  \textbf{Metastable start}: crossing+3000 (Paper 7 convention).
\item
  \textbf{\(q_1\ge q_2\ge q_3\ge q_4\)}: singular values from the
  dominant-plane Gram reduction. \textbf{rank\_q} is the effective rank
  under the existing relative threshold. rank\_q is the rank of
  \(Q=[B_0\mid B_{\mathrm{dom}}]\) (an \(M\times4\) matrix) and is
  structurally capped at 4.
\item
  \textbf{Five-component occupation decomposition}: occupations of
  directions 1--4 (two parent-plane directions plus two additional
  directions), the residual rotating subspace, and the kernel. The
  occupations sum to 1 at every step (projection closure error 0).
\item
  \textbf{Direction subspace} \(D_{34}(t)\): the orthonormal basis of
  the two additional directions returned by the existing
  \texttt{s4\_new\_dirs}. The canonical comparison object is not the
  columns but the projector
\end{itemize}

\[P_{34}(t)=D_{34}(t)D_{34}(t)^{\mathsf T},\qquad \mathcal O(P,Q)=\tfrac12\operatorname{tr}(PQ),\]

with \(\mathcal O=1\) meaning subspace identity. Individual direction
3/4 columns depend on basis choice (column exchange, sign) and are not
used for lineage judgment.

\begin{itemize}
\tightlist
\item
  \textbf{First-passage coordinates}
  \(\tau_k=\min\{t\mid f(t)\ge 10^{-k}\}\): first passage of each
  decimal level. The interval amplification rate is defined as
  \(\ln 10/\Delta\tau\). These are coordinates of the amplification
  process, not event thresholds.
\end{itemize}

\subsubsection{1.2 The two explicit seeds}\label{the-two-explicit-seeds}

The Paper 7 orbit contains two artificial inputs.

\begin{enumerate}
\def\labelenumi{\arabic{enumi}.}
\tightlist
\item
  \textbf{Initial seed}: a zero-closure pair in the parent kernel added
  at \(\delta=10^{-15}\) (\(f(0)\approx\delta^2=10^{-30}\)).
\item
  \textbf{Metastable seed}: a single transverse perturbation
  \(\varepsilon\eta_\perp\) (\(\varepsilon=10^{-8}\)) injected once at
  \(t_1=\mathrm{crossing}+3000\) (the perturbation-response measurement
  of Paper 7 Figure 4).
\end{enumerate}

This paper removes the two independently and separates which observed
facts causally depend on the seeds.

\subsubsection{1.3 Structure of the initial state: two-directionality
fixed by zero-square closure
{[}Result{]}}\label{structure-of-the-initial-state-two-directionality-fixed-by-zero-square-closure-result}

For \(Z=X+iY\),

\[Z^{\mathsf T}Z=\lVert X\rVert^2-\lVert Y\rVert^2+2i\,X\cdot Y,\]

so the zero-square closure \(Z^{\mathsf T}Z=0\) is equivalent to

\[\lVert X\rVert=\lVert Y\rVert,\qquad X\cdot Y=0.\]

Every nonzero state satisfying the closure condition necessarily fixes a
pair of equal-norm orthogonal directions in the real \(M\)-dimensional
relational space. The two-directional structure is not a choice of
initial condition but a consequence of the first axiom.

Measurements of the \texttt{make\_parent} output are given in Table 1.

\emph{Table 1: Structure of the parent state (originals verified by
SHA-256; same PRNG as Stage A2a)}

{\def\LTcaptype{none} % do not increment counter
\begin{longtable}[]{@{}
  >{\raggedright\arraybackslash}p{(\linewidth - 4\tabcolsep) * \real{0.3333}}
  >{\raggedright\arraybackslash}p{(\linewidth - 4\tabcolsep) * \real{0.3333}}
  >{\raggedright\arraybackslash}p{(\linewidth - 4\tabcolsep) * \real{0.3333}}@{}}
\toprule\noalign{}
\begin{minipage}[b]{\linewidth}\raggedright
Quantity
\end{minipage} & \begin{minipage}[b]{\linewidth}\raggedright
N=5 (M=10)
\end{minipage} & \begin{minipage}[b]{\linewidth}\raggedright
N=40 (M=780)
\end{minipage} \\
\midrule\noalign{}
\endhead
\bottomrule\noalign{}
\endlastfoot
Nonzero complex components / M & 10 / 10 & 780 / 780 \\
min \(\lvert v_e\rvert\) & 0.2887 & 0.0324 \\
rank\([\mathrm{Re}\,v,\ \mathrm{Im}\,v]\) & \textbf{2} & \textbf{2} \\
\(\lVert\mathrm{Re}\,v\rVert=\lVert\mathrm{Im}\,v\rVert\) & \(1/\sqrt2\)
(error \(10^{-15}\)) & \(1/\sqrt2\) (error \(10^{-15}\)) \\
\(\mathrm{Re}\,v\cdot\mathrm{Im}\,v\) & \(-4.4\times10^{-16}\) &
\(-3.6\times10^{-16}\) \\
\(\lvert v^{\mathsf T}v\rvert\) & \(1.1\times10^{-15}\) &
\(2.3\times10^{-15}\) \\
Invariant-plane residual \(\lVert KX-\mu Y\rVert\) &
\(2.0\times10^{-13}\) & \(1.1\times10^{-13}\) \\
Decomposition dimensions of the initial \(K\) (parent / rotating
complement / kernel) & 2 / 4 / 4 & 2 / 78 / 700 \\
\end{longtable}
}

Table 1 establishes three facts. First, all \(M\) relational waves are
nonzero; this is not a state keeping only two particular components.
Second, the real rank of the whole state is 2: all relational waves are
collapsed onto a single two-dimensional rotation plane (an invariant
plane of the generator \(K\), residual \(10^{-13}\)). Third,
\texttt{make\_parent} imposes only the closure condition and
self-consistency (being an eigenmode of the \(K\) built from its own
phases); neither a phase-quantization rule nor the coordinates of two
axes are given from outside.

The separate examination referred to in the Claims --- that the closure
condition inherently contains quantization --- is the expository note
{[}3{]}. From the two requirements of zero-square-sum constraint and
scale anonymity, the state space becomes a projective quadric in complex
projective space (a compact Kähler manifold), and discrete spectra,
integer Chern classes, and the qubit isomorphism at \(N=3\) follow from
citations of known mathematics alone. The implementation here does not
impose that quantization rule from outside; the two-direction collapse
in Table 1 is a measurement of the double-precision dynamical system.

Whether the initial collapse onto two directions is the direct cause of
three-direction formation is not identified (§9).

\begin{center}\rule{0.5\linewidth}{0.5pt}\end{center}

\subsection{2. Experimental design}\label{experimental-design}

\subsubsection{2.1 Conditions and experimental
systems}\label{conditions-and-experimental-systems}

{\def\LTcaptype{none} % do not increment counter
\begin{longtable}[]{@{}
  >{\raggedright\arraybackslash}p{(\linewidth - 6\tabcolsep) * \real{0.2143}}
  >{\centering\arraybackslash}p{(\linewidth - 6\tabcolsep) * \real{0.2857}}
  >{\centering\arraybackslash}p{(\linewidth - 6\tabcolsep) * \real{0.2857}}
  >{\raggedright\arraybackslash}p{(\linewidth - 6\tabcolsep) * \real{0.2143}}@{}}
\toprule\noalign{}
\begin{minipage}[b]{\linewidth}\raggedright
Condition
\end{minipage} & \begin{minipage}[b]{\linewidth}\centering
Initial seed
\end{minipage} & \begin{minipage}[b]{\linewidth}\centering
Metastable seed
\end{minipage} & \begin{minipage}[b]{\linewidth}\raggedright
Initial state
\end{minipage} \\
\midrule\noalign{}
\endhead
\bottomrule\noalign{}
\endlastfoot
A (fully seedless) & OFF & OFF & \(Z_0=v\) (no random numbers
consumed) \\
B (initial only) & ON & OFF & \(Z_0=(v+\delta g)/\lVert\cdot\rVert\) \\
D (Paper 7 equivalent) & ON & ON & same as B; single injection of
\(\varepsilon\eta_\perp\) at \(t_1\) \\
\end{longtable}
}

\begin{itemize}
\tightlist
\item
  \textbf{Preliminary experiment} (N=5, 40, 300, t \(\le\) 55000): full
  comparison of conditions A/B/D.
\item
  \textbf{Stage A2a} (N=5, every-step recording, t \(\le\) 5000):
  high-density observation of the fully seedless orbit, starting from
  the numerical floor \(f(0)=3.275\times10^{-33}\) of \(Z_0=v\).
\item
  \textbf{Stage A2c / A2d} (N=5): direction-lineage tracking of
  \(P_{34}(t)\) for the seeded / fully seedless orbits.
\item
  \textbf{Long-horizon control} (N=5, 40, 300, fully seedless, t \(\le\)
  110000): search for a second rapid expansion.
\item
  \textbf{Transition anatomy} (N=5, Paper 7 reproduced orbit): temporal
  separation of rank\_q, occupations, and f.
\end{itemize}

\subsubsection{2.2 Hypotheses tested}\label{hypotheses-tested}

The hypotheses remaining at the close of Paper 7 were fixed in advance.

\begin{itemize}
\tightlist
\item
  \textbf{H-lower-bottom}: a lower metastable phase (a second bottom)
  with an isomorphic growth-and-arrest process is hidden inside the flat
  region before the rapid expansion.
\item
  \textbf{H-seed-causation}: the explicit seeds determine the
  occurrence, the timing, or the first generated direction of the
  splitting.
\item
  \textbf{H-scaled-copy}: the small direction subspace before the rapid
  expansion is a scaled copy of the late three directions.
\item
  \textbf{H-repeating-cycle}: the metastable state after the rapid
  expansion is an intermediate phase preparing the next rapid expansion.
\end{itemize}

The results reject or fail to support all four, and in doing so
establish the time structure stated in the Claims.

\begin{center}\rule{0.5\linewidth}{0.5pt}\end{center}

\subsection{3. Result 1: the full sequence reproduces with the seeds
removed
{[}Result{]}}\label{result-1-the-full-sequence-reproduces-with-the-seeds-removed-result}

Condition A (fully seedless) reproduced crossing, rank\_q increase, and
metastabilization at every N (Table 2).

\emph{Table 2: Conditions A/B/D at all N (t \(\le\) 55000)}

{\def\LTcaptype{none} % do not increment counter
\begin{longtable}[]{@{}
  >{\raggedleft\arraybackslash}p{(\linewidth - 16\tabcolsep) * \real{0.1111}}
  >{\raggedright\arraybackslash}p{(\linewidth - 16\tabcolsep) * \real{0.1111}}
  >{\raggedleft\arraybackslash}p{(\linewidth - 16\tabcolsep) * \real{0.1111}}
  >{\raggedleft\arraybackslash}p{(\linewidth - 16\tabcolsep) * \real{0.1111}}
  >{\raggedleft\arraybackslash}p{(\linewidth - 16\tabcolsep) * \real{0.1111}}
  >{\raggedleft\arraybackslash}p{(\linewidth - 16\tabcolsep) * \real{0.1111}}
  >{\raggedleft\arraybackslash}p{(\linewidth - 16\tabcolsep) * \real{0.1111}}
  >{\raggedleft\arraybackslash}p{(\linewidth - 16\tabcolsep) * \real{0.1111}}
  >{\raggedleft\arraybackslash}p{(\linewidth - 16\tabcolsep) * \real{0.1111}}@{}}
\toprule\noalign{}
\begin{minipage}[b]{\linewidth}\raggedleft
N
\end{minipage} & \begin{minipage}[b]{\linewidth}\raggedright
Cond.
\end{minipage} & \begin{minipage}[b]{\linewidth}\raggedleft
crossing
\end{minipage} & \begin{minipage}[b]{\linewidth}\raggedleft
metastable start
\end{minipage} & \begin{minipage}[b]{\linewidth}\raggedleft
max\_f
\end{minipage} & \begin{minipage}[b]{\linewidth}\raggedleft
final\_f
\end{minipage} & \begin{minipage}[b]{\linewidth}\raggedleft
final rank\_Q
\end{minipage} & \begin{minipage}[b]{\linewidth}\raggedleft
mean\_meta(f)
\end{minipage} & \begin{minipage}[b]{\linewidth}\raggedleft
std\_meta(f)
\end{minipage} \\
\midrule\noalign{}
\endhead
\bottomrule\noalign{}
\endlastfoot
5 & A & 1166 & 4166 & 0.9655 & 0.8754 & 4 & 0.8744 &
\(1.29\times10^{-2}\) \\
5 & B & 1167 & 4167 & 0.9600 & 0.8053 & 4 & 0.8054 &
\(2.87\times10^{-2}\) \\
5 & D & 1167 & 4167 & 0.9600 & 0.8026 & 4 & 0.8034 &
\(2.88\times10^{-2}\) \\
40 & A & 2011 & 5011 & 0.1938 & 0.1938 & 4 & 0.1934 &
\(7.38\times10^{-4}\) \\
40 & B & 2011 & 5011 & 0.2035 & 0.2035 & 4 & 0.2030 &
\(1.29\times10^{-3}\) \\
40 & D & 2011 & 5011 & 0.2035 & 0.2035 & 4 & 0.2030 &
\(1.29\times10^{-3}\) \\
300 & A & 4849 & 7849 & 0.0898 & 0.0898 & 4 & 0.0882 &
\(3.53\times10^{-3}\) \\
300 & B & 4844 & 7844 & 0.0861 & 0.0861 & 4 & 0.0848 &
\(2.94\times10^{-3}\) \\
300 & D & 4844 & 7844 & 0.0861 & 0.0861 & 4 & 0.0848 &
\(2.94\times10^{-3}\) \\
\end{longtable}
}

\begin{figure}
\centering
\pandocbounded{\includegraphics[keepaspectratio,alt={f comparison N=5}]{figs/fig01_f_compare_N00005.png}}
\caption{f comparison N=5}
\end{figure}

\emph{Figure 1: f(t) for conditions A/B/D at N=5. The crossing
difference between seedless (A) and seeded (B/D) is 1 step. N=40 and 300
are isomorphic.}

The crossing differences are 1 step at N=5, 0 steps at N=40, and 5 steps
at N=300. The explicit initial seed contributes neither to the
occurrence nor to the timing of the splitting (first part of
H-seed-causation rejected).

\begin{center}\rule{0.5\linewidth}{0.5pt}\end{center}

\subsection{4. Result 2: no separate bottom inside the latent region
{[}Result{]}}\label{result-2-no-separate-bottom-inside-the-latent-region-result}

The fully seedless N=5 orbit (Stage A2a, every-step recording) was
inspected over its entire structure from the initial numerical floor
\(f(0)=3.275\times10^{-33}\) to the main growth.

\subsubsection{4.1 Sign inspection of every-step
differences}\label{sign-inspection-of-every-step-differences}

For all 1129 differences from step 0 to just before the first passage of
\(f\ge10^{-2}\) (step 1129):

\begin{itemize}
\tightlist
\item
  positive differences: \textbf{1129 (all)}
\item
  negative differences: \textbf{0}
\item
  zero differences: \textbf{0}
\end{itemize}

The shape ``growth \(\to\) arrest \(\to\) lower metastable shelf \(\to\)
regrowth'' exists nowhere in the every-step data from the initial floor
to the main growth. The flatness in linear display is an appearance
caused by the smallness of the observable; in logarithmic display it is
one continuous rise (Figure 2).

\begin{figure}
\centering
\pandocbounded{\includegraphics[keepaspectratio,alt={seedless log10 f}]{figs/figure02_seedless_log10_f.png}}
\caption{seedless log10 f}
\end{figure}

\emph{Figure 2: \(\log_{10}f\) of the fully seedless N=5 orbit. A single
continuous rise from the numerical floor.}

\subsubsection{4.2 First-passage series and amplification
rate}\label{first-passage-series-and-amplification-rate}

The passage steps from \(10^{-30}\) to \(10^{-23}\) rise as 1, 2, 5, 14,
36, 68, 107, 150, and over all 21 intervals of \(10^{-23}\to10^{-2}\)
the passage time converges to \textbf{46 or 47 steps per decade}. The
median interval exponential rate is \(4.95\times10^{-2}\) (natural log
per step).

Linear regression of \(\ln f\) over the 980 points from \(f\ge10^{-23}\)
(step 150) to \(f\ge10^{-2}\) (step 1129) gives

\[\lambda=0.04937,\qquad R^2=0.9999993.\]

The per-step ratio of f is \(e^{\lambda}=1.0506\); the amplitude ratio
is \(e^{\lambda/2}=1.0250\). The main part of the latent region is
described by a single geometric law (Figure 3).

\begin{figure}
\centering
\pandocbounded{\includegraphics[keepaspectratio,alt={decade rate comparison}]{figs/figure05_decade_growth_rate_comparison.png}}
\caption{decade rate comparison}
\end{figure}

\emph{Figure 3: Mean exponential rate per decade. Seeded and seedless
coincide at every level.}

\subsubsection{4.3 Judgment}\label{judgment}

\textbf{H-lower-bottom is rejected.} The scope of the rejection is N=5,
double precision, every-step recording (t \(\le\) 5000). Structure below
the double-precision floor, and unperturbed departure in exact
arithmetic, are outside the scope of this paper (§9).

\begin{center}\rule{0.5\linewidth}{0.5pt}\end{center}

\subsection{5. Result 3: the seeds select neither timing nor initial
direction
{[}Result{]}}\label{result-3-the-seeds-select-neither-timing-nor-initial-direction-result}

\subsubsection{5.1 Timing}\label{timing}

Comparing first passages of the seeded and seedless orbits at all
levels, the first passage of \(f\ge10^{-12}\) is absolute step 662 for
both, and the per-decade passage-time difference is at most 1 step over
all intervals (Figure 4).

\begin{figure}
\centering
\pandocbounded{\includegraphics[keepaspectratio,alt={absolute step comparison}]{figs/figure03_seeded_vs_seedless_absolute_step.png}}
\caption{absolute step comparison}
\end{figure}

\emph{Figure 4: f(t) of seeded and seedless orbits overlaid on absolute
steps. They coincide without any translation.}

The initial values are inverted. The seeded orbit starts from
\(f(0)=1.066\times10^{-30}\) and the seedless orbit from
\(f(0)=3.275\times10^{-33}\), and the two merge onto the same geometric
series within a few hundred steps. The explicit seed \(\delta=10^{-15}\)
is smaller than the eigenmode residual \(2.140\times10^{-13}\) of the
double-precision parent. Even in the seeded orbit, what governed the
effective minimal transverse component was not the explicit seed.

\subsubsection{5.2 Initial direction}\label{initial-direction}

The same-step direct comparison of \(P_{34}(t)\) between the seedless
and seeded orbits is given in Table 3 (Stage A2d).

\emph{Table 3: Same-step overlap, seedless vs.~seeded}

{\def\LTcaptype{none} % do not increment counter
\begin{longtable}[]{@{}lrrr@{}}
\toprule\noalign{}
Interval & Samples & Overlap median & Overlap minimum \\
\midrule\noalign{}
\endhead
\bottomrule\noalign{}
\endlastfoot
Pre-crossing, resolvable & 40 & \textbf{0.999999987} & 0.999976 \\
crossing--1799 & 128 & 0.999997 & 0.994306 \\
1800--2500 & 141 & 0.999852 & 0.990155 \\
after 2500 & 500 & 0.265980 & 0.203573 \\
\end{longtable}
}

From before the rapid expansion until just after the establishment of
the three-direction closure, the direction subspaces of the two orbits
are identical. \textbf{The explicit initial seed does not select the
first generated direction subspace} (second part of H-seed-causation
rejected).

The divergence after step 2500 (overlap first below 0.95 at step 2690,
below 0.5 at step 3185) is a long-time orbit difference separated from
the initial selection of the generated direction. After this divergence,
rank\_q=4 and the occupation levels are retained in both orbits (§6.3).

\begin{center}\rule{0.5\linewidth}{0.5pt}\end{center}

\subsection{6. Result 4: rank\_q=4 onset, occupation growth, and the f
expansion occur at different times
{[}Result{]}}\label{result-4-rank_q4-onset-occupation-growth-and-the-f-expansion-occur-at-different-times-result}

\subsubsection{6.1 rank\_q=4 appears at the numerical
floor}\label{rank_q4-appears-at-the-numerical-floor}

\begin{itemize}
\tightlist
\item
  Seedless orbit: the first rank\_q=4 is at saved step \textbf{120},
  where \(f=2.089\times10^{-24}\); the bracketing saved occupation rows
  (steps 107 / 125) have direction 3/4 occupations at the
  \(10^{-30}\)--\(10^{-29}\) level.
\item
  Seeded orbit: the first saved rank\_q=4 is at step \textbf{265},
  \(f=3.127\times10^{-21}\), with bracketing direction 3/4 occupations
  at \(4\times10^{-28}\)--\(8\times10^{-28}\). The number of rank\_q=4
  saved rows before crossing is 97.
\end{itemize}

rank\_q=4 appears while splitting fraction and direction occupations are
at the numerical floor. Therefore

\[\mathrm{rank_q}=4\quad\not\Rightarrow\quad\text{establishment of a direction structure with finite occupation},\]

and the onset of rank\_q=4 cannot serve as the time of three-direction
establishment (Figure 5).

\begin{figure}
\centering
\pandocbounded{\includegraphics[keepaspectratio,alt={q ratios and rank\_q}]{figs/figure06_q_ratios_and_rank_q_0_3000.png}}
\caption{q ratios and rank\_q}
\end{figure}

\emph{Figure 5: \(q_3/q_1\), \(q_4/q_1\) and rank\_q at N=5 (steps
0--3000). The early rank\_q response and the occupation growth occur at
different times.}

\subsubsection{6.2 Quantities that grow with delays after
crossing}\label{quantities-that-grow-with-delays-after-crossing}

After crossing=1167, direction 3/4 occupations, the kernel, and the q
ratios do not change at the same speed (Table 4).

\emph{Table 4: Observables after crossing (N=5, actual saved values)}

{\def\LTcaptype{none} % do not increment counter
\begin{longtable}[]{@{}rrrrrr@{}}
\toprule\noalign{}
step & f & d3 occ. & d4 occ. & \(q_3/q_1\) & \(q_4/q_1\) \\
\midrule\noalign{}
\endhead
\bottomrule\noalign{}
\endlastfoot
1200 & 0.133 & 0.0027 & 0.0035 & 0.279 & 0.016 \\
1400 & 0.585 & 0.023 & 0.195 & 0.608 & 0.248 \\
1800 & 0.840 & 0.213 & 0.210 & 0.733 & 0.507 \\
2500 & 0.801 & 0.319 & 0.045 & 0.752 & 0.549 \\
\end{longtable}
}

The f crossing, the finite growth of direction occupations, and the
growth of the q ratios are separated over hundreds of steps; a
single-event description ``the three directions were established at one
time'' is not supported by the data. That direction 3/4 occupations grow
to finite magnitudes (the \(10^{-1}\) level) after the rapid expansion
and are retained is an observational fact, as in Table 4 and Figure 5.

\subsubsection{6.3 What is metastable is the occupation
structure}\label{what-is-metastable-is-the-occupation-structure}

Conditions A and B diverge as orbits in the metastable region (§5.2),
while rank\_Q=4, the ceiling-level f, and the finiteness of the
direction occupations are preserved in both (Table 2). The
three-direction closure is not three fixed axes; it is a dynamical
closure that preserves the number of occupied dimensions and the
occupation distribution while its internal directions can mix over long
times.

\begin{center}\rule{0.5\linewidth}{0.5pt}\end{center}

\subsection{7. Result 5: the direction subspace reorganizes during the
rapid expansion
{[}Result{]}}\label{result-5-the-direction-subspace-reorganizes-during-the-rapid-expansion-result}

\subsubsection{7.1 The pre-expansion direction is not a scaled copy of
the late
direction}\label{the-pre-expansion-direction-is-not-a-scaled-copy-of-the-late-direction}

In the fully seedless orbit (Stage A2d), the overlap between the
\(P_{34}\) that became resolvable before the rapid expansion
(resolvability criterion \(\min(q_3,q_4)/q_1\ge10^{-6}\); 40 samples at
steps 985--1165) and the late representative subspace (top two
eigenvectors of the mean projector over steps 1800--2500) is

\[\mathcal O_{\mathrm{early,late}}\ \text{median}=0.1425\qquad(\text{range}\ [0.1210,\ 0.1445]).\]

The seeded orbit (Stage A2c) gives the same level (median 0.1495). A
quantity that would approach 1 for identical subspaces stays at 0.14.
\textbf{The small direction subspace before the rapid expansion is not a
scaled copy of the late three directions} (H-scaled-copy rejected;
Figure 6).

\begin{figure}
\centering
\pandocbounded{\includegraphics[keepaspectratio,alt={seedless direction lineage}]{figs/figure01_seedless_early_vs_late_lineage.png}}
\caption{seedless direction lineage}
\end{figure}

\emph{Figure 6: Early-vs-late direction lineage of the fully seedless
N=5 orbit. The overlap between the early \(P_{34}\) and the late
representative stays at 0.14.}

\subsubsection{7.2 The reorganization is a real rotation during the
expansion}\label{the-reorganization-is-a-real-rotation-during-the-expansion}

In the fixed observation window steps 900--1400 of the seeded orbit, the
sum of maximal principal angles between consecutive-step \(P_{34}\)
(rotation path length) is 3.871 rad, and the maximum of the
consecutive-step maximal principal angle is 0.372 rad. Column exchanges
are 0; this is not a relabeling of basis columns but rotation and mixing
of the subspace itself. Stages A2c/A2d classify this lineage as
\textbf{ROTATING\_OR\_MIXED\_LINEAGE}.

\subsubsection{7.3 Reclassification of Paper 7 Figure
4}\label{reclassification-of-paper-7-figure-4}

The overlap between the transverse perturbation directions \(T_\perp\)
applied in the metastable region in Paper 7 Figure 4 and the late
representative direction is, for the three perturbation seeds,

\[\mathcal O(D_{34}^{\mathrm{late}},T_\perp)=0.0330,\ 0.0301,\ 0.0182.\]

The artificially excitable response channels coincide neither with the
naturally occupied three directions nor with the precursor directions
before the rapid expansion. What Paper 7 Figure 4 measured was not
``sprouts of additional directions arising naturally'' but
\textbf{response channels outside the three-direction closure that are
excitable but not naturally occupied}.

\begin{center}\rule{0.5\linewidth}{0.5pt}\end{center}

\subsection{8. Result 6: the metastable state does not proceed to a
second rapid expansion
{[}Result{]}}\label{result-6-the-metastable-state-does-not-proceed-to-a-second-rapid-expansion-result}

\subsubsection{8.1 Removal of the second seed: the metastable
oscillation is not a product of the
seed}\label{removal-of-the-second-seed-the-metastable-oscillation-is-not-a-product-of-the-seed}

Under condition B (no metastable seed), the temporal variation in the
metastable region continued to the final step (std\_meta:
\(2.87\times10^{-2}\) at N=5, \(1.29\times10^{-3}\) at N=40,
\(2.94\times10^{-3}\) at N=300). The difference from condition D (single
injection of \(\varepsilon=10^{-8}\) at \(t_1\)) is exact agreement
within recording precision at N=40 and 300, and about \(10^{-3}\) at N=5
(Figure 7).

\begin{figure}
\centering
\pandocbounded{\includegraphics[keepaspectratio,alt={metastable B vs D}]{figs/fig04_metastable_B_vs_D_N00005.png}}
\caption{metastable B vs D}
\end{figure}

\emph{Figure 7: Conditions B and D in the metastable region at N=5. The
transverse perturbation does not change the metastable oscillation.}

The metastable oscillation is not a product of the
transverse-perturbation seed; it is a consequence of the closed
dynamics. After the injection in condition D, no second crossing or rank
increase appears within the observation range t \(\le\) 55000.

\subsubsection{8.2 Long-horizon control: no sign of a second rapid
expansion up to
t=110000}\label{long-horizon-control-no-sign-of-a-second-rapid-expansion-up-to-t110000}

The fully seedless orbits were extended to t=110000 at N=5, 40, 300.

\begin{itemize}
\tightlist
\item
  N=40, 300: the raw f, \(q_3\), \(q_4\) over t=55000--110000 form
  \textbf{no extremum at all} and approach their terminal values
  monotonically. The peak-to-peak over t=90000--110000 is
  \(f=1.0\times10^{-11}\) at N=40 and \(f=2.1\times10^{-11}\) at N=300
  (recording-precision level).
\item
  N=5: the metastable oscillation (f range {[}0.807, 0.966{]} over steps
  4166--55000) decays by t=55000, and over the extension 55000--110000 f
  is constant at 0.875392 (peak-to-peak \(1.0\times10^{-11}\),
  recording-precision floor). In the extension, steps exceeding the
  metastable maximum: 0; drops below f=0.05: 0; second crossings: 0.
\item
  The late relaxation cannot be absorbed by a single exponential; a
  double exponential removes the residual valley (AIC improvements at
  N=300: f=2848.6, \(q_3\)=1568.8, \(q_4\)=2172.0; Figure 8).
\item
  The early metastable region contains raw-value reversals of a damped
  oscillation (N=40: \(q_3\) maximum at t=1650, minimum at t=1700).
  N=300 is recorded at 100-step intervals, leaving shorter oscillations
  possibly unresolved.
\end{itemize}

\begin{figure}
\centering
\pandocbounded{\includegraphics[keepaspectratio,alt={long-horizon residual}]{figs/figure_long_horizon_one_exp_residual_x20.png}}
\caption{long-horizon residual}
\end{figure}

\emph{Figure 8: Single-exponential residual of the long-horizon control.
The remaining structure is explained by double-exponential relaxation;
there is no reversal toward re-amplification.}

Within the observation range \textbf{H-repeating-cycle is not
supported}. The metastable state is a long-lived phase that, after the
initial damped oscillation, relaxes multi-exponentially toward a steady
level; it contains no growth series heading toward a second geometric
expansion.

\subsubsection{8.3 No precursor of a bosonic/fermionic distinction
appears in the observation
range}\label{no-precursor-of-a-bosonicfermionic-distinction-appears-in-the-observation-range}

This series has confirmed in separate experimental systems that wave
interactions have two types. In contrast to the linear interaction in
which waves exchange and mix while remaining superposable, the nonlinear
exchange-interference scattering with reflection and transmission
exchanges localization between an extended wave and a localized wave
{[}4{]} and forms metastable intermediate states (the
white-cat/black-cat/gray-cat interface) {[}5{]}. In those experiments
the scattering rule is given from outside by the experimenter as the
interaction type.

The experimental system of this paper has no direct observable that
discriminates this type. The update is a single unitary step on the
whole state by the state-dependent generator; it has no test direction
constructing a per-wave-pair reflection/transmission scattering process.
What follows is therefore not a designed test but a passive precursor
search over the existing observables (time series, direction
occupations, rank, metastabilization process).

The operational content of the search is whether an additional discrete
direction beyond three, or a new discrete channel, becomes naturally
occupied. If the separation of linear and nonlinearly scattering waves
arises internally along with the rapid expansion or direction formation,
its first trace appears as natural occupation of additional directions.
The result is Table 5.

\emph{Table 5: Residual rotating occupation in the metastable region of
the fully seedless natural orbits (occupation outside directions 1--4
and the kernel)}

{\def\LTcaptype{none} % do not increment counter
\begin{longtable}[]{@{}
  >{\raggedleft\arraybackslash}p{(\linewidth - 6\tabcolsep) * \real{0.2500}}
  >{\raggedleft\arraybackslash}p{(\linewidth - 6\tabcolsep) * \real{0.2500}}
  >{\raggedright\arraybackslash}p{(\linewidth - 6\tabcolsep) * \real{0.2500}}
  >{\raggedright\arraybackslash}p{(\linewidth - 6\tabcolsep) * \real{0.2500}}@{}}
\toprule\noalign{}
\begin{minipage}[b]{\linewidth}\raggedleft
N
\end{minipage} & \begin{minipage}[b]{\linewidth}\raggedleft
Dimension of residual rotating subspace
\end{minipage} & \begin{minipage}[b]{\linewidth}\raggedright
Metastable max
\end{minipage} & \begin{minipage}[b]{\linewidth}\raggedright
Terminal value
\end{minipage} \\
\midrule\noalign{}
\endhead
\bottomrule\noalign{}
\endlastfoot
5 & 2 & \(1.5\times10^{-3}\) & \(4.4\times10^{-16}\) \\
40 & 76 & \(4.7\times10^{-6}\) & 0 \\
300 & 596 & \(9.7\times10^{-8}\) & \(2.1\times10^{-17}\) \\
\end{longtable}
}

Since rank\_q is structurally capped at 4 (§1.1) and cannot detect
additional directions, the judgment rests on this direct measurement.
The residual rotating occupation grows at no N and decays to numerical
zero (Figure 9). The response channels outside the closure are excitable
but not naturally occupied (§7.3). The long-horizon control (t \(\le\)
110000) likewise shows no growth series of a new channel.

\begin{figure}
\centering
\pandocbounded{\includegraphics[keepaspectratio,alt={five-component occupation}]{figs/figure3_compare_N5_N40_N300.png}}
\caption{five-component occupation}
\end{figure}

\emph{Figure 9: Five-component occupation decomposition of the fully
seedless natural orbits (N=5, 40, 300). No direction outside the three
directions and the kernel becomes naturally occupied.}

No sign by which a spontaneous separation into bosonic and fermionic
states could be judged exists in these observables. That this separation
causes the rapid expansion or three-direction formation is not excluded,
and not supported. What is confirmed is that the transition from the
two-directional state to the metastable three-directional structure
occurs without installing a fermionic scattering rule.

\begin{center}\rule{0.5\linewidth}{0.5pt}\end{center}

\subsection{9. Boundary conditions (what this paper does not
prove)}\label{boundary-conditions-what-this-paper-does-not-prove}

\begin{enumerate}
\def\labelenumi{\arabic{enumi}.}
\tightlist
\item
  \textbf{Double-precision floor}: the minimal input of the seedless
  orbit is the eigenmode residual of the parent state
  (\(2.14\times10^{-13}\)) and rounding error. Unperturbed departure in
  exact arithmetic --- the physical origin of the first minimal
  difference --- is not separated here. This is not a condition
  weakening the central claims but the next object to be separated.
\item
  \textbf{N and time range}: the every-step inspection for the
  lower-bottom rejection (§4) is N=5, t \(\le\) 5000. The long-horizon
  control (§8.2) is N=5, 40, 300, t \(\le\) 110000. Repetition with
  still longer periods is not logically excluded, but no raw-data
  reversal indicating such a candidate exists in the present range.
\item
  \textbf{Recording intervals}: q is saved every 5 steps and occupations
  every 25--100 steps; simultaneity claims are limited to the actually
  saved records. Sample-to-sample rotation quantities in the lineage
  analysis are not one-step maxima.
\item
  \textbf{Incomplete control}: the numerical-resolution sweep (control
  with an artificial quantization operator) has incomplete aggregation.
  The completed part is a control in which cutting off the minimal
  amplification source at finite resolution pins the orbit to the parent
  state; it is not evidence of a second bottom.
\item
  \textbf{Correspondence with physical XYZ}: the correspondence between
  the formed three-direction closure and physical spatial directions is
  a separate connection problem handled with projection Gram matrices,
  principal angles, and interference quantities. What this series
  establishes is that three directions became readable from the \(M\)
  phase relations, not that they are the xyz of a background space. The
  physical canonical readout is the interference of direction planes
  (inner products of normals).
\item
  \textbf{Causality between the initial two directions and
  three-direction formation}: that the initial state collapses onto two
  directions as a consequence of the first axiom (§1.3) is a measured
  fact, but whether this collapse directly causes three-direction
  formation is not identified. A systematic control from generic
  closure-only states (\texttt{zero\_closure\_generic} initialization
  without the eigenmode condition) is the next task.
\item
  \textbf{Inherent phase quantization}: the viewpoint of the expository
  note {[}3{]} that the closure condition inherently contains
  quantization is not imposed as an external rule in this implementation
  and is not verified by the present data.
\item
  \textbf{Untestability of the statistical distinction}: the precursor
  search of §8.3 is a passive search constrained by the absence of a
  direct observable discriminating the interaction type; it is not an
  experiment testing spontaneous formation of the statistical
  distinction. Testing it requires designing an observation system that
  constructs scattering processes inside the closed relational system
  (generation of localized wave packets and readout of collisions) --- a
  connection problem between the exchange-interference scattering
  systems {[}4,5{]} and the present system.
\end{enumerate}

\begin{center}\rule{0.5\linewidth}{0.5pt}\end{center}

\subsection{10. Reproducibility}\label{reproducibility}

\begin{itemize}
\tightlist
\item
  The dynamical code reuses the Paper 7 originals
  (\texttt{run\_n\_scaling\_lowrank\_v1.py} and others) fixed by SHA-256
  and imported read-only; the original folders are unchanged. All hashes
  are recorded in the stage reports and
  \texttt{config/source\_file\_hashes.json}.
\item
  The Paper 7 N=5 orbit was reproduced bitwise in both the f and q
  series (Stage A0). The seedless orbit is bitwise identical across two
  independent executions (Stage A2a exec 1/2). The N=300 five-component
  decomposition was cross-checked against the existing condition A with
  all rows of the 27 common columns identical.
\item
  Numerical health: normalization error \(\le1.4\times10^{-14}\);
  zero-square closure \(\le1.8\times10^{-10}\) (immediately after the
  condition-D injection; otherwise \(\le1.3\times10^{-13}\)); projection
  closure error 0; occupation-sum error 0.
\item
  Canonical sources of the numbers: Table 1, §4.1, the §4.2 regression,
  §6.1 seedless, §8.2 N=5, and Table 5 are in
  \texttt{paper8\_draft\_supporting\_analysis\_v1/reports/}. Table 2 and
  §8.1 are in \texttt{reports/preliminary\_seed\_ablation\_report.md}.
  §4.2 and §5.1 are in
  \texttt{paper8\_stage\_A2a\_seedless\_N5/reports/}. Table 3 and §7.1
  are in
  \texttt{paper8\_stage\_A2d\_seedless\_direction\_lineage\_N5/reports/}.
  §6 is in \texttt{paper7\_N5\_transition\_anatomy/reports/}. §7.2 and
  §7.3 are in
  \texttt{paper8\_stage\_A2c\_direction\_lineage\_N5/reports/}. §8.2 is
  in
  \texttt{paper7\_seedless\_natural\_figures3\_4\_v1/validation\_report.md}.
\item
  Environment: Python 3.9.6, NumPy 2.0.2, macOS arm64. Random numbers
  via \texttt{numpy.random.default\_rng} with fixed seeds.
\end{itemize}

\begin{center}\rule{0.5\linewidth}{0.5pt}\end{center}

\subsection{11. Conclusion}\label{conclusion}

The phenomenon discovered in Paper 7 --- the long low-change region, the
geometric rapid expansion, and the spontaneous formation of the
metastable three-directional structure --- reproduces at all of N=5, 40,
300 with both explicit seeds removed. The seeds cause none of the
occurrence, the timing, the first generated direction, or the metastable
oscillation.

The initial two-directional state at the starting point is not two axes
placed by the experimenter. Zero-square closure forces equal-norm
orthogonal two-directionality on every nonzero state, and all \(M\)
relational waves of \texttt{make\_parent}, while nonzero, collapse
spontaneously onto a single rotation plane of real rank 2. The observed
sequence of this series begins with a two-direction closure
spontaneously established from the axioms and ends with the dynamical
formation of a three-direction closure.

The seemingly flat region before the rapid expansion is a single
continuous rise with zero negative differences under every-step
inspection --- a long latent process with no lower bottom inside. The
early onset of rank\_q=4 is a response at the numerical floor and cannot
be identified with direction establishment. The direction subspace
rotates, mixes, and reorganizes during the rapid expansion; the three
directions are dynamically selected and formed inside the expansion
process, not laid down in advance. The formed metastable state relaxes
multi-exponentially after the initial damped oscillation and produces no
second rapid expansion up to the observation range t=110000.

The open problems are narrowed to three. (i) Why does the geometric
growth begin after the long latency --- including the physical origin of
the first minimal difference below the double-precision floor. (ii) Why
does the reorganization during the expansion select three directions.
(iii) Why does the formed metastable state settle down instead of
proceeding to the next expansion. Direct verifications are: for (i),
reproduction of the latent region in exact or high-precision arithmetic;
for (ii), comparison of the most unstable eigenplane of the transverse
Jacobian in co-rotating coordinates (the 144-step Floquet matrix) with
the measured precursor \(P_{34}\); for (iii), identification of the
relaxation spectrum on the metastable shell.

The possibility that the separation of bosonic and fermionic
interactions underlies the phenomenon remains neither excluded nor
supported (§8.3). Three-direction formation without an installed
fermionic scattering rule is confirmed. Whether the statistical
distinction of waves appears afterwards as a result of this formation
process, or requires a degree of freedom not observable in this system,
is a task for the next experimental series, together with the design of
an observation system constructing scattering processes inside the
closed relational system.

\subsubsection{11.1 Relation to prior
work}\label{relation-to-prior-work}

\begin{itemize}
\tightlist
\item
  \textbf{Latency followed by geometric growth} shares its time
  structure with modulational instability, in which a uniform carrier
  destabilizes exponentially into transverse sidebands from a noise
  floor {[}8{]}. In the present system, however, what grows is not
  spatial sidebands but the direction subspace itself, and the
  directions reorganize during amplification (§7).
\item
  The \textbf{transverse spectrum of a state-dependent generator} is a
  toolset shared with the locked-state spectrum analysis of coupled
  phase oscillators {[}7{]}. The present system, however, carries the
  norm-structural constraint of zero-square closure, and what is read
  out is direction structure, not synchronization.
\item
  \textbf{Long-time metastabilization without a second expansion} (§8)
  can be read in the vocabulary of the FPUT metastable state {[}6{]} and
  prethermalization in isolated systems {[}9{]} --- a long-lived
  nonthermal intermediate state. No recurrence corresponding to the FPUT
  type appears within t \(\le\) 110000.
\item
  The prior work closest in claim shape --- three directions alone
  expanding out of a symmetric candidate space --- is the dynamical
  emergence of (3+1)-dimensional spacetime in the IIB (IKKT) matrix
  model {[}10{]}, where three directions extend via spontaneous breaking
  of SO(9,1)/SO(10) symmetry. Mechanism and level both differ: against a
  path integral over matrix degrees of freedom, this paper is a
  deterministic closed dynamics of N-body relational waves, and the
  three directions are not dimensions of a background space but
  direction structure read out from the interference of relational
  waves. Whether the shared selection of ``3'' is another expression of
  one mechanism is left as an open question.
\end{itemize}

No prior work corresponds to the central results of this paper --- the
causal removal of explicit seeds, the rejection of a lower bottom by
every-step continuity of the latent region, and the direct measurement
of direction-subspace reorganization during the rapid expansion.

\begin{center}\rule{0.5\linewidth}{0.5pt}\end{center}

\subsection{References}\label{references}

{[}1{]} N. Kihara, Basic Axiom System of the Anonymous Equal-Amplitude
Composite Wave Model (purified definition paper), Zenodo. Concept DOI:
\href{https://doi.org/10.5281/zenodo.21315735}{10.5281/zenodo.21315735}.

{[}2{]} N. Kihara, Arrest of Spontaneous Splitting and Creation of
Additional Axes in Closed N-Body Relational Wave Systems (Paper 7),
Zenodo. Concept DOI:
\href{https://doi.org/10.5281/zenodo.21543070}{10.5281/zenodo.21543070}.

{[}3{]} N. Kihara, The Geometric Identity of the Zero-Square-Sum
Constraint and Scale Anonymity --- Projectivized Isotropic Cone,
Projective Quadric, and Inherent Quantumness (expository note), Zenodo.
Concept DOI:
\href{https://doi.org/10.5281/zenodo.21495305}{10.5281/zenodo.21495305}.

{[}4{]} N. Kihara, How Is Localization Exchanged by the
Exchange-Interference Scattering Matrix?, Zenodo. Concept DOI:
\href{https://doi.org/10.5281/zenodo.21333766}{10.5281/zenodo.21333766}.

{[}5{]} N. Kihara, Reading the White-Cat/Black-Cat/Gray-Cat Metastable
States in a Closed System, Zenodo. Concept DOI:
\href{https://doi.org/10.5281/zenodo.21353208}{10.5281/zenodo.21353208}.

{[}6{]} G. Benettin et al., The Metastable State of
Fermi--Pasta--Ulam--Tsingou Models, Entropy \textbf{25}(2), 300 (2023).

{[}7{]} R. E. Mirollo, S. H. Strogatz, The spectrum of the locked state
for the Kuramoto model of coupled oscillators, Physica D \textbf{205},
249 (2005).

{[}8{]} I. Daumont, T. Dauxois, M. Peyrard, Modulational instability:
first step towards energy localization in nonlinear lattices,
Nonlinearity \textbf{10}, 617 (1997).

{[}9{]} T. Mori, T. N. Ikeda, E. Kaminishi, M. Ueda, Thermalization and
prethermalization in isolated quantum systems, J. Phys. B \textbf{51},
112001 (2018).

{[}10{]} S.-W. Kim, J. Nishimura, A. Tsuchiya, Expanding
(3+1)-dimensional universe from a Lorentzian matrix model for
superstring theory in (9+1)-dimensions, Phys. Rev.~Lett. \textbf{108},
011601 (2012).

\end{document}
